\chapter{\texttt{defpyfun}: Funciones Python con Captura Lisp}
\label{chap:defpyfun}
{\huge$\bigstar$}{\huge$\bigstar$}{\huge$\bigstar$} Avanzado

\section{Definición}
\begin{lispcode}
(let ((factor 3.1416))
  (defpyfun scale (x) "return x * factor" :as scale-py))
(scale-py 10)
\end{lispcode}

\section{Mecanismo de Captura}
Las variables libres se serializan como argumentos por defecto:
\begin{verbatim}
def scale(x, _captured_factor=3.1416):
    factor = _captured_factor
    return x * factor
\end{verbatim}

\begin{examplebox}[title=Leaky ReLU con parámetros capturados]
\begin{lispcode}
(let ((alpha 0.01))
  (defpyfun leaky-relu (x) "return x if x > 0 else alpha * x" :as lrelu-py))
(lrelu-py -2)
\end{lispcode}
\end{examplebox}

